\documentclass[11pt]{article}
\usepackage[margin=1in]{geometry}
\usepackage[T1]{fontenc}
\usepackage{amsmath,amssymb}
\usepackage[expansion=false]{microtype}
\usepackage{enumitem}
\usepackage{xcolor}
\usepackage{listings}
\usepackage[hidelinks]{hyperref}

\definecolor{codebg}{rgb}{0.965,0.965,0.965}
\definecolor{codecmt}{rgb}{0.33,0.5,0.33}
\definecolor{codekw}{rgb}{0.20,0.20,0.6}
\lstset{
  basicstyle=\ttfamily\footnotesize,
  backgroundcolor=\color{codebg},
  commentstyle=\color{codecmt}\itshape,
  keywordstyle=\color{codekw}\bfseries,
  breaklines=true,columns=fullflexible,keepspaces=true,
  showstringspaces=false,frame=single,framesep=4pt,xleftmargin=4pt,
  morekeywords={def,returns,raises,asserts,for,if,else,yield,class,interface,artifact}
}
\newcommand{\artifact}[1]{\textbf{\texttt{#1}}}
\setlist{topsep=2pt}

\title{\textbf{Thought Anchors --- Implementation \& Execution Specification}}
\author{Companion build spec to the research protocol\\ \emph{(thought\_anchors\_experiments\_v3)}}
\date{}

\begin{document}
\maketitle
\vspace{-2em}

\section{Purpose and how to use this document}
This is a \emph{build specification}, not source code. It states what code \emph{should} exist---module
responsibilities, data schemas, interface signatures, algorithms in sketch form, smoke-test assertions,
the parallel/sequential phase structure, the cluster run model, the per-result decision logic, and the
human breakpoints---at a level explicit enough to implement directly, without committing the
implementations yet. It covers everything in the protocol \emph{through the trials/tiers}; the causal
interventions (Q1--Q4) are out of scope here and built only after a detector is validated.

\paragraph{The two-layer build philosophy (read first).} The document is written all at once so that
early plumbing anticipates later plumbing; the \emph{code} is built in two layers:
\begin{itemize}
\item \textbf{Foundation (\S\ref{sec:foundation}): build fully, up front, against forward-looking
interfaces.} The shared spine---schema, storage, checkpointing, splits, raw-rollout retention,
config assertion, smoke harness. This is where ``no leak / no mid-run crash / launch-and-forget'' is
won or lost, so it is designed with the \emph{whole} pipeline in view, including signals not computed
until later tiers.
\item \textbf{Tier modules (\S\ref{sec:modules}, \S\ref{sec:phases}): specify now, generate and run on
demand at the gates.} The tiered/fork design exists precisely so machinery is added only as evidence
demands. Build a whole tier, run it, decide at the breakpoint; do not work ahead.
\end{itemize}
``Wait for results'' means wait at the \textbf{breakpoints} (\S\ref{sec:breakpoints}), not after every
module.

\section{Repository layout}
\begin{lstlisting}
ta/
  config/        run configs (yaml): model, hyperparams, phase, tier, seeds
  core/
    modelio.py   M0  load + config assertion
    traces.py    M1  base CoT trace generation
    spans.py     M2  clause-respecting span segmentation
    signals/     M3  one file per signal family (force, katz, mlp_g, grad_d, online)
    labels.py    M4  resample-and-filter gold labels (+ constructed-swap check)
    combiners.py M5  ridge / PCR / MLP
    metrics.py   M6  spearman, bootstrap CI, receiver-head baseline, decision logic
  infra/
    schema.py        dataclasses + (de)serialization for all artifacts
    store.py         atomic keyed read/write, content/config hashing
    queue.py         file-based work queue, claim/complete sentinels
    seeds.py         deterministic seed derivation
    splits.py        grouped-by-problem split assignment + leak assertions
    manifest.py      run manifest (git sha, config hash, model facts)
  smoke/         one test module per core module + integration tests
  run/
    launch_phase.sh  tmux launcher: one worker pane per GPU + monitor
    worker.py        idempotent unit processor (pulls from queue)
    summarize.py     end-of-phase decision report + DONE sentinel
  data/          raw datasets + adapters (babi, bbh_shuffled, clutrr, scone, gsm8k)
  artifacts/     all outputs, sharded by phase/dataset/problem_id (gitignored)
\end{lstlisting}

\section{Environment and dependencies}
\begin{itemize}
\item Python 3.11; \texttt{torch} (CUDA build matching the cluster driver), \texttt{transformers},
\texttt{accelerate}, \texttt{safetensors}, \texttt{numpy}, \texttt{scipy} (Spearman, Fisher-z),
\texttt{scikit-learn} (ridge/PCR/MLP, grouped CV), \texttt{pyarrow} (parquet), \texttt{pyyaml}.
\item Two model handles (\S\ref{sec:m0}): an \emph{eager}-attention handle for the attribution pass
(needed to materialise attention probabilities) and an SDPA/flash handle for the resampling rollouts.
\item Pin all versions in a lockfile; record the lockfile hash in the run manifest. Transformers
internals (attention return signature, RoPE, \texttt{repeat\_kv}) are version-sensitive; assert against
them in M0 rather than trusting the version string.
\end{itemize}

\section{Foundation: cross-cutting design (build this first, completely)}
\label{sec:foundation}
Everything here is shared infrastructure with high change-cost. Implement and smoke-test it before any
signal or label code, against interfaces that already include the not-yet-built signals.

\subsection{Artifact schema (\texttt{infra/schema.py})}
Forward-looking: the \artifact{Signals} record reserves fields for \emph{every} signal, even those a
given run does not populate, so later tiers never reshape storage.
\begin{lstlisting}
Problem:   problem_id, dataset, phase, prompt, gold_answer,
           vocabulary: {type -> [valid values]},   # typed closed vocab (structured tasks)
           supporting_facts: [token-spans]|None,     # bAbI free ground truth
           meta
Trace:     trace_id, problem_id, token_ids, text, gen_seed, gen_config_hash
Span:      span_id, trace_id, clause_id, tok_start, tok_end, text,
           is_window: bool, task_tokens: [idx], frame_tokens: [idx]
Signals:   span_id ->
             E, P, Pchained,          # routing (one-hop, Katz)
             G,                        # computation (MLP write magnitude)
             Pprime, PprimeChained,    # crude computation-weighted outflow
             D, Dchained, D_ig,        # gradient-scaled outflow (+ chained, + IG)
             c, position, length,      # online, confounds
             meta: {band, K, lambda, g, sigma_pop}
Label:     span_id ->
             A,                        # gold importance (Eq. label)
             real_dist: {answer->p}, diff_dist: {answer->p},
             raw_rollouts: [{cond: real|diff, answer, parse_ok, filter_sim}],  # RETAINED
             R_real, R_diff_net, R_gross, parse_fail_rate, metric: tv|kl|flip
ConstructReport: span_id -> {diff_shift, similar_shift, has_task_tokens}
SplitAssignment: problem_id -> {fold, role: train|test}
RunManifest: git_sha, config_hash, lockfile_hash, model_facts(M0), seeds, phase, tier, started_at
\end{lstlisting}
\textbf{Critical no-leak decision:} \artifact{Label.raw\_rollouts} is retained in full. The expensive
artifact is the rollouts, not $A$; storing per-rollout answers means $A$ (and the metric choice, and the
filter threshold) can be recomputed offline forever \emph{without re-rolling}. Never store only the
scalar $A$.

\subsection{Storage (\texttt{infra/store.py})}
\begin{itemize}
\item Sharded by \texttt{phase/dataset/problem\_id}; tabular as parquet/jsonl, tensors as safetensors.
\item \textbf{Atomic writes}: write to \texttt{*.tmp} then \texttt{rename} (POSIX atomic); a half-written
file can never be read as complete.
\item \textbf{Content/config keying}: each artifact's path embeds a hash of its inputs + config. A rerun
with identical inputs/config is a no-op (idempotent); a changed hyperparameter writes a new key rather
than silently overwriting.
\item \texttt{read(key)} raises on absent (callers distinguish ``not computed'' from ``empty'').
\end{itemize}

\subsection{Checkpointing, idempotency, resumability (\texttt{infra/queue.py})}
\begin{itemize}
\item Unit of work = the smallest re-runnable chunk: one \emph{problem} for trace gen / signal
extraction; one \emph{span} (or small span-batch) for rollouts (rollouts dominate cost, so checkpoint
finely there).
\item File-based queue: \texttt{claim(unit)} writes a \texttt{.claimed} lease (with PID + timestamp);
\texttt{complete(unit)} writes a \texttt{.done} sentinel. On restart a worker skips \texttt{.done} units
and reclaims stale \texttt{.claimed} leases (PID dead or lease older than $T$).
\item Consequence: kill the job at any point and relaunch; it resumes from the last \texttt{.done} with
no recompute and no corruption. This is the launch-and-forget guarantee, and it is smoke-tested
(\S\ref{sec:smoke}, resume-after-kill).
\end{itemize}

\subsection{Determinism and seeding (\texttt{infra/seeds.py})}
\begin{lstlisting}
seed(purpose, problem_id) -> int   # = hash(global_seed, purpose, problem_id) mod 2^31
\end{lstlisting}
All generation/sampling draws its seed this way, so any unit reproduces identically on rerun. Record the
global seed in the manifest. Note throughput-vs-bitwise-determinism: set deterministic kernels only for
the smoke tests that assert bitwise equality; production rollouts may relax this for speed but remain
\emph{seeded} (distributions reproduce, not necessarily bitwise).

\subsection{Grouped splits and leak prevention (\texttt{infra/splits.py})}
\begin{itemize}
\item Splits are assigned at \emph{problem} granularity; a problem's every span shares its problem's
fold/role. Within-trace spans are correlated, so splitting them across train/test leaks.
\item \texttt{assert\_no\_leak(split)}: train and test \texttt{problem\_id} sets are disjoint; fails the
run loudly if violated. Called before any combiner fit and inside the combiner smoke test.
\item Effective sample size for all CIs is the \emph{number of traces/problems}, not spans.
\end{itemize}

\subsection{Config assertion as a hard gate (\texttt{infra/manifest.py} + M0)}
The architecture corrections in the protocol all hinge on the exact model config; assert them, do not
trust them. See \S\ref{sec:m0}.

\subsection{Smoke-test harness (\texttt{smoke/})}
A single runner discovers and runs all smoke tests, prints a pass/fail table, and \textbf{blocks phase
launch} unless the phase's required tests pass. Every module ships its smoke tests with it. The harness
distinguishes \emph{correctness} tests (must pass to proceed) from \emph{calibration} reports (inform
breakpoints).

\section{Module specifications}
\label{sec:modules}
Each module lists: responsibility, interface, key correctness points (from the protocol), and smoke
tests. Interfaces are signatures, not implementations.

\subsection{M0 --- Model loading and config assertion (\texttt{core/modelio.py})}
\label{sec:m0}
\textbf{Responsibility.} Load the two model handles; extract and assert architecture facts; expose the
static weights the signal code needs ($W_O$ per-head blocks, RMSNorm gains).
\begin{lstlisting}
load_model(path, attn_impl: "eager"|"sdpa") -> (model, tokenizer)
assert_config(model.config) -> ModelFacts
  asserts: num_hidden_layers == 28
           hidden_size == 1536 ; head_dim == 128
           num_attention_heads (H) and num_key_value_heads (Hkv) present
           group_size = H // Hkv  (record; GQA expected, Hkv < H)
           norm type is RMSNorm ; record rms_norm_eps
           rope present ; record theta
           eager handle actually returns attention probs on a 4-token probe
ModelFacts: {L,H,Hkv,group_size,d,d_head,eps,...}  # written to manifest
wo_head_blocks(model) -> Tensor[L, H, d_head, d]   # o_proj reshaped per query head
rms_gains(model) -> {input_ln[L], post_attn_ln[L]}  # gamma vectors
\end{lstlisting}
\textbf{Correctness.} Fail loudly on any mismatch --- each is a silent-corruption risk (GQA mis-grouping,
wrong norm linearisation, attention probs silently empty under SDPA). This runs as step 1 of every phase.

\textbf{Smoke tests.} (i) assert passes on the real model; (ii) assert \emph{raises} on a synthetic
config with wrong head counts / norm type; (iii) eager probe returns attention of shape $[H,T,T]$;
(iv) \texttt{wo\_head\_blocks} recomposed equals the original \texttt{o\_proj} weight.

\subsection{M1 --- Base trace generation (\texttt{core/traces.py})}
\textbf{Responsibility.} Generate the CoT traces that are the substrate for spans, signals, and labels.
\begin{lstlisting}
generate_trace(problem, model_fast, gen_config, seed) -> Trace
  # few-shot prompt forcing a parseable final answer ("Answer: X")
  # temperature, max_new_tokens, stop conditions in gen_config
\end{lstlisting}
\textbf{Correctness.} Deterministic per \texttt{seed(problem\_id,"trace")}; store full token\_ids (spans
index into these). Record the answer-format scaffold per dataset (M4 parses it).

\textbf{Smoke tests.} 5 bAbI problems generate, parse to a valid answer, and round-trip through
store/load unchanged; rerun with same seed is bitwise identical.

\subsection{M2 --- Clause-respecting span segmentation (\texttt{core/spans.py})}
\textbf{Responsibility.} Two-stage boundary-first segmentation (protocol \S Granularity): parse clauses,
then window within a clause.
\begin{lstlisting}
segment(trace, tokenizer, t, clause_rules) -> [Span]
  # stage 1: split decoded text at sentence terminators, newlines,
  #          comma clause-boundaries, connectives {so,therefore,thus,then,next,because}
  # stage 2: map each clause's char span -> token span via offset_mapping
  #          (tokenizer(..., return_offsets_mapping=True));
  #          clause <= t tokens -> 1 span; longer -> consecutive t-token windows
  #          (final remainder merged into previous if < t/2). Never cross a clause boundary.
tag_vocab(span, problem.vocabulary) -> (task_tokens, frame_tokens)  # for construct check
\end{lstlisting}
\textbf{Correctness.} The text$\leftrightarrow$token alignment is the subtle part --- use offset mappings,
never re-tokenise substrings. Spans partition the trace (every non-special token in exactly one span).

\textbf{Smoke tests.} (i) spans never cross a clause boundary; (ii) decoded span tokens equal the clause
substring they came from (alignment); (iii) union of spans = full trace, disjoint; (iv) windows respect
$t$ and the merge rule.

\subsection{M3 --- Signal extraction, the attribution pass (\texttt{core/signals/})}
\textbf{Responsibility.} One eager forward pass over the base trace, plus per-layer local backward passes,
producing all forward-time signals at span granularity.

\paragraph{Delivered force and edges (\texttt{force.py}).}
\begin{lstlisting}
delivered_force_norms(trace, model_eager, facts) -> Ehat[T,T], (per-target topK survivors)
  # per query head h: F^{l,h}_{s->t} = WO_block[h] @ (alpha^{l,h}_{t,s} * v^{kv(h)}_s)
  #   GQA: kv(h) = h // group_size ; value from shared kv head
  #   pre-norm: F added to residual with NO post-attn norm -> exactly additive, no bias
  # pruning: Ehat_{s->t} = sum_{l,h} alpha * ||WO_block[h] @ v_s||  (cheap upper bound)
  #   keep top-K sources per target t -> N(t); compute exact ||F|| only on survivors
edges(...) -> E[span,span]   # E_{s->t} = sum_l ||F^l_{s->t}||, aggregated to spans
\end{lstlisting}
Correctness: GQA mapping; $W_O$ per-head block; do \emph{not} apply any post-attention norm (pre-norm).
Pruning is top-$K$ incoming per target by the routing bound $\widehat E$; with bounded $g$ this also
covers the one-hop computation signals (use bounded $g$, below).

\paragraph{Katz routing propagation (\texttt{katz.py}).}
\begin{lstlisting}
katz(E, lambda) -> P[span]   # right-to-left DP on the causal DAG:
                             # P_s = sum_{t>s} E_{s->t} (1 + lambda * P_t)
\end{lstlisting}
DAG (strict $t>s$) guarantees termination; sweep $\lambda$ later (Tier 3). One-hop $E$-outflow is
$\lambda{=}0$.

\paragraph{Computation magnitude (\texttt{mlp\_g.py}).}
\begin{lstlisting}
mlp_write_G(trace, model_eager, band B, sigma_pop) -> G[span]
  # per layer l in B: g_l = || MLP_l( RMSNorm_postattn(h'_l) ) ||   (down_proj output, SwiGLU)
  # G_t = sum_{l in B} g_l / sigma_l ;  G_span = sum_{t in span} G_t
  # sigma_l = cross-position std at layer l, computed over sigma_pop (DECIDE: pooled vs per-trace),
  #           treated as a frozen constant (detached)
\end{lstlisting}
Correctness: in pre-norm $\lVert\text{MLP output}\rVert$ \emph{is} the net residual write --- do not
compute a separate ``net write''. Band $B$ localised in Phase 0.

\paragraph{Gradient-scaled attribution (\texttt{grad\_d.py}).}
\begin{lstlisting}
grad_D(trace, model_eager, band B, sigma) -> D[span], (optional) D_ig[span]
  # per layer l in B (LOCAL Jacobian, not full-network gradient):
  #   treat pre-MLP residual h'_l as a detached leaf requiring grad
  #   forward through layer-l MLP sub-block (incl. its input RMSNorm)
  #   loss = sum_t || MLP_l ||  (sum over positions; position-wise => no cross-contamination)
  #   backward once -> grad_{h'_l}[t] = d||MLP_l||/d h'_{l,t}   for all t
  #   D^l_{s->t} = (grad_{h'_l}[t] / sigma_l) . F^l_{s->t}
  # D_{s->t} = sum_{l in B} D^l ;  aggregate to spans like E
  # D_ig: replace point gradient by integral over alpha in [0,1] of grad at h'_t - (1-alpha) F_{s->t}
  #       (baseline h'_t - F_{s->t}, shared with Taylor and exact ablation)
\end{lstlisting}
Correctness: differentiate w.r.t.\ the \emph{residual} $h'$ (RMSNorm in the autodiff path) so it composes
with $F$; one backward \emph{per layer} in $B$, position-vectorised; $\sigma$ detached. Cost $\approx|B|$
local backwards.

\paragraph{Computation-weighted outflow and chained variants (\texttt{outflow.py}).}
\begin{lstlisting}
pprime(E, G, g) -> Pprime[span]            # P'_i = sum_{j>i} E_{i->j} g(G_j)
chained(E_or_D, G, g, lambda) -> [span]    # Tier-3: Katz-style recursion carrying g(G)
online_c(N(t)) -> c[span]                  # streaming in-degree (# targets that kept i)
position, length -> from spans
\end{lstlisting}
Correctness: restrict $g$ to \emph{bounded} forms (log / centred / threshold) so routing-pruning suffices
and chained variants do not blow up; chaining is Tier-3 (improvement) or a Tier-2 rescue, not Tier-1.

\textbf{M3 smoke tests (the correctness core --- gate the whole project).}
\begin{itemize}
\item \textbf{Force additivity (the key test).} Reconstruct each target's pre-MLP residual from the
per-source forces: $h_{\text{block-in}} + \sum_s F_{s\to t} \stackrel{?}{\approx} h'_t$ (actual residual),
within fp tolerance. Passing validates GQA grouping, $W_O$ blocks, $\alpha$ capture, and pre-norm
additivity \emph{simultaneously}.
\item \textbf{GQA mapping.} Per-head $F$ on a tiny input matches a hand-computed reference.
\item \textbf{Gradient finite-difference.} $\nabla G\cdot\delta \approx G(h'{+}\delta)-G(h')$ for small
random $\delta$ (validates the local Jacobian).
\item \textbf{D vs exact ablation} on a span subsample: $D_{s\to t}$ correlates with
$G_t(h'_t)-G_t(h'_t-F_{s\to t})$ --- so a future $D$-null is read as ``no signal'' not ``bad
approximation''.
\item \textbf{Katz vs brute force} on a tiny hand-built DAG.
\item \textbf{Pruning soundness}: $\widehat E \ge \lVert F\rVert$ elementwise (valid upper bound).
\end{itemize}

\subsection{M4 --- Resampling labels (\texttt{core/labels.py})}
\textbf{Responsibility.} Produce the gold importance $A$ by the resample-and-filter method of the
protocol; provide the constructed-swap construct-validity check; retain raw rollouts.
\begin{lstlisting}
resample_filter(trace, span, model_fast, R, filter, metric, seed) -> Label
  # real:  K_real continuations of the unperturbed prefix-through-i
  # diff:  resample replacement spans for i FROM THE MODEL (prefix up to start of i),
  #        filter for semantically-DIFFERENT (filter.sim(orig,repl) < tau), keep until R_diff net;
  #        continue CoT to answer from each
  # parse each continuation -> discrete answer (per-dataset parser); drop unparseable (record rate)
  # A = metric( real_dist, diff_dist )   # TV default; match theirs (tv|kl|flip)
  # store ALL raw_rollouts (cond, answer, parse_ok, filter_sim)
filter: SemanticFilter   # FROM BOGDAN ET AL. CODE if reused; else reimplement (embed/NLI + tau)
parse_answer(text, dataset) -> answer|None
constructed_swap_check(trace, span, problem.vocabulary) -> ConstructReport
  # different = swap a task-content token to same-type sibling; similar = edit frame tokens only
  # report diff_shift (should be high on anchors) and similar_shift (should be ~0)
\end{lstlisting}
\textbf{Correctness / decisions.}
\begin{itemize}
\item \textbf{Reuse first (Tier-0).} If Bogdan et al.'s released labels cover the model, ingest them and
skip generation entirely --- this fixes the filter, the metric, and the hyperparameters, and sidesteps
the granularity question. Only reimplement if reuse fails.
\item \textbf{Filter / metric / hyperparameters} (\texttt{tau}, oversampling factor to net $R$ different,
$K_{\text{real}}$, temperature, max-continuation) are read from their code; do not invent.
\item \textbf{Granularity transfer.} Their filter is sentence-tuned; in Phase 0 confirm it calibrates on
clause-windows (semantically-different resamples of a short span are obtainable). If not, fall back to
sentence granularity (which aligns with reusing their dataset).
\item \textbf{Rollouts on the SDPA/flash handle} (fast); the eager handle is only for M3.
\end{itemize}
\textbf{M4 smoke tests.}
\begin{itemize}
\item \textbf{Real-vs-real null (the key test).} With no perturbation, $A\approx0$ --- establishes the
label noise floor so any $A>0$ is signal, not pipeline artifact.
\item \textbf{Constructed-swap behaviour} on bAbI: \texttt{diff\_shift} large on annotated supporting
facts, \texttt{similar\_shift}$\approx0$; a real-vs-real null per arm.
\item \textbf{Parser accuracy} $\ge$ threshold on a hand-labelled sample; \texttt{parse\_fail\_rate}
recorded; require a minimum count of valid rollouts per condition or discard the span.
\item \textbf{Raw-rollout recompute}: recomputing $A$ from stored \texttt{raw\_rollouts} reproduces the
stored $A$ (so labels never need regeneration).
\item \textbf{Filter sanity}: known paraphrase pairs classified similar, known content-swaps different.
\end{itemize}

\subsection{M5 --- Combiners (\texttt{core/combiners.py}) [Tier 2+]}
\textbf{Responsibility.} Combine signals into a predictor; the interpretable model's coefficients are
the deliverable, the MLP is an accuracy ceiling.
\begin{lstlisting}
fit_ridge(signals, labels, split) -> {coef, rho_test, ci}        # alpha by grouped CV
fit_pcr(signals, labels, split) -> {...}                          # SVD principal-component regression
fit_mlp(signals, labels, split) -> {rho_test, ci}                 # 2-layer, width 32-64, early stop
\end{lstlisting}
\textbf{Correctness.} \texttt{assert\_no\_leak} on the split before every fit; grouped K-fold by problem;
report on held-out problems only; keep the outflow signal only if it adds over $P{+}G$.

\textbf{Smoke tests.} grouped split shares no problem across train/test; ridge recovers known
coefficients on synthetic data; reruns reproduce.

\subsection{M6 --- Metrics, baseline, decision logic (\texttt{core/metrics.py})}
\textbf{Responsibility.} Score signals against $A$, reproduce the receiver-head baseline, apply the
two-pronged decision.
\begin{lstlisting}
spearman_vs_A(signal, labels, by_trace) -> (rho, ci)   # bootstrap over TRACES; Fisher-z
receiver_head_baseline(trace, model_eager) -> scores    # reuse Bogdan et al. code if available
tier1_decision(rhos) -> {prong_i: any rho>0.22?,
                         prong_ii: max(rho_G, rho_Pprime) > rho_P (or adds over P)?,
                         fork: improvement|fix}
\end{lstlisting}
\textbf{Correctness.} Effective $N$ = number of traces; pre-declare the confirmatory comparisons and the
multiple-comparison correction (FDR over the marginals) separately from exploratory sweeps.

\textbf{Smoke tests.} Spearman matches \texttt{scipy} on known data; bootstrap CI has nominal coverage on
synthetic data with known $\rho$; baseline reproduces a known value on a fixture trace.

\section{Smoke-test suite (consolidated)}
\label{sec:smoke}
Correctness tests block launch; calibration reports inform breakpoints. Beyond the per-module tests
above, the integration tests:
\begin{itemize}
\item \textbf{End-to-end dry run} on 5 problems: traces $\to$ spans $\to$ signals $\to$ small-$R$ labels
$\to$ one ridge fit, all artifacts written and reloadable, no exception.
\item \textbf{Resume-after-kill}: launch the dry run, \texttt{SIGKILL} a worker mid-rollout, relaunch;
assert (a) no \texttt{.done} unit recomputed, (b) no corrupt/half-written artifact, (c) final result
identical to an uninterrupted run.
\item \textbf{Leak audit}: across the whole pipeline, no \texttt{problem\_id} appears in two folds; no
span's label rollout reuses its own continuation as a ``different'' sample.
\item \textbf{Determinism}: two seeded dry runs produce identical artifact hashes (under deterministic
kernels).
\item \textbf{Resource guards}: pre-flight estimates artifact bytes and peak GPU memory from
(\#problems, trace length, $R$, $T$); asserts free disk and fits batch to memory; refuses launch
otherwise.
\item \textbf{Schema round-trip}: every artifact type serialises and deserialises to an equal object,
including the reserved-but-empty signal fields.
\end{itemize}
\textbf{Rule:} a phase launches only when its required correctness tests are green. This is what makes
``launch and walk away'' safe.

\section{Phased execution and the cluster run model}
\label{sec:phases}
Tiers (\emph{what} to run) and curriculum phases (\emph{which} data) are orthogonal. Tier 0/1 run on the
Phase-0/1 data first; later tiers extend to later phases. Within any cluster run:

\paragraph{Parallel vs sequential.}
\begin{itemize}
\item \textbf{Embarrassingly parallel} (one queue unit each, fan out across all 8 GPUs): trace generation
(per problem), signal extraction (per trace), and---dominant---resampling rollouts (per span). These
never need to talk to each other.
\item \textbf{Sequential barriers}: segmentation after traces; signals and labels may run concurrently
(both depend only on traces); analysis after both; band/$t$/$K$ localisation needs a \emph{first small
label batch}, so a short sequential pre-batch precedes the big parallel labeling.
\end{itemize}

\paragraph{The tmux launch-and-forget model (\texttt{run/}).}
\begin{lstlisting}
launch_phase.sh <phase> <tier>:
  1. run required smoke tests; abort if any correctness test fails
  2. build the work queue for the phase (list of unit ids; skip .done)
  3. start a tmux session:
       - one pane per GPU: worker.py --gpu k   (CUDA_VISIBLE_DEVICES=k; pull/claim/process/complete)
       - one monitor pane: progress, ETA, failure tail, queue depth
  4. on empty queue: summarize.py writes the decision report + a DONE sentinel
worker.py: idempotent; on crash, lease expires and another worker reclaims the unit; checkpoints per unit
\end{lstlisting}
You launch, detach tmux, and leave. Workers are idempotent and resumable, artifacts are atomic and keyed,
so a node failure or preemption costs at most the in-flight units, never a re-run. The DONE sentinel +
report tell you when to come back.

\paragraph{Run-phase mapping.}
\begin{enumerate}
\item \textbf{Build phase (local, sequential).} Implement the Foundation; pass all correctness smoke
tests; end-to-end dry run. $\Rightarrow$ \textbf{Breakpoint 1}.
\item \textbf{Tier-0 gates (cheap, mostly local).} Reuse check; free bAbI construct check; small-$R$
pilot \emph{that doubles as} $t$/$K$/band-$B$ localisation and the granularity-transfer check.
$\Rightarrow$ \textbf{Breakpoint 2}.
\item \textbf{Tier-1 run (cluster, parallel).} Full $R$ labeling + $P,G,P'$ marginals + receiver-head
baseline on Phase-0/1. $\Rightarrow$ \textbf{Breakpoint 3 (the fork)}.
\item \textbf{Tier-2/3 (cluster, on demand).} Improvement (or fix) features per the fork; combiners;
chaining/reach; extend phases. $\Rightarrow$ \textbf{Breakpoint 4}.
\end{enumerate}

\section{Decision logic --- what to do per result}
\label{sec:decisions}
Mirrors the protocol's tiers; computed by \texttt{tier1\_decision} and the gate scripts.
\begin{itemize}
\item \textbf{Tier-0 reuse.} Covered $\Rightarrow$ ingest their labels, skip generation, proceed to
Tier-1 marginals. Not covered $\Rightarrow$ reimplement filter/metric from their code.
\item \textbf{Tier-0 construct check.} Signals rank supporting facts above distractors $\Rightarrow$
encouraging, proceed. They do not $\Rightarrow$ strong warning; inspect signal code before spending the
labeling budget (leading indicator, not a hard stop).
\item \textbf{Tier-0 pilot.} Even noisy $\rho$ shows nothing on any signal $\Rightarrow$ stop and
re-examine before full $R$. Granularity-transfer fails $\Rightarrow$ fall back to sentence granularity
(and prefer reuse).
\item \textbf{Tier-1 two-pronged.} (i) some marginal $>0.22$? (ii) does $G$ or $P'$ beat / add over $P$?
\begin{itemize}
\item both yes $\Rightarrow$ \emph{improvement}: add $D$ (Trial 3) and the linear combiner; keep outflow
only if it adds over $P{+}G$.
\item routing clears but computation does not (or position-dominated) $\Rightarrow$ \emph{fix}: rule out
(a) position confound (residualise), (b) bad approximation ($D^{\mathrm{IG}}$), (c) wrong reach (pull the
chained variant forward as a rescue) before declaring the axis inert.
\item nothing clears $0.22$ $\Rightarrow$ fix branch, then if still null, a clean negative result.
\end{itemize}
\item \textbf{Tier-3.} Chained underperforms or position-dominated $\Rightarrow$ remediation ladder
(centre $g$ $\to$ per-hop cap $\to$ sub-critical $\lambda$ $\to$ row-normalise; residualise position).
MLP ceiling reported beside, never instead of, the interpretable model.
\end{itemize}

\section{Human breakpoints}
\label{sec:breakpoints}
Stop, inspect code and results, and decide --- do not let the pipeline auto-continue past these.
\begin{enumerate}
\item \textbf{After the Foundation (before any expensive labeling).} Review the correctness core:
force-additivity, gradient finite-difference, D-vs-ablation, real-vs-real null. \emph{Question:} is the
machinery provably correct? Nothing costly runs until yes.
\item \textbf{After Tier-0 gates.} Review reuse availability, the construct-check ranking, the pilot
$\rho$, and the chosen $t/K/B$ and granularity-transfer verdict. \emph{Question:} are inputs and
hyperparameters sane, and is the full labeling budget worth committing?
\item \textbf{After Tier-1 marginals (the fork).} Review the two-pronged decision. \emph{Question:} did
the fulcrum clear $0.22$, and did the \emph{computation axis} pay off (vs.\ a routing-only win)? Choose
the improvement or fix branch here.
\item \textbf{After the first combiner / chained results.} \emph{Question:} is the gain real, or
position/overfit? Confirm on held-out problems and against the position-only null before scaling phases.
\end{enumerate}

\section{Build-and-run checklist}
\begin{enumerate}
\item Implement \texttt{infra/} (schema, store, queue, seeds, splits, manifest) + smoke harness.
\item Implement M0; pass config-assertion tests.
\item Implement M1, M2, M3; pass the correctness core (force additivity, gradient FD, Katz, pruning).
\item Implement M4 (with constructed-swap check) + M6 baseline/metrics; pass real-vs-real null and parser
tests. \emph{End-to-end dry run.} $\rightarrow$ \textbf{Breakpoint 1}.
\item Tier-0: reuse check $\to$ construct check $\to$ small-$R$ pilot/localisation. $\rightarrow$
\textbf{Breakpoint 2}.
\item Launch Tier-1 on the cluster (tmux). $\rightarrow$ \textbf{Breakpoint 3}.
\item Implement M5 and the fork-selected features; launch Tier-2/3 on demand. $\rightarrow$
\textbf{Breakpoint 4}.
\end{enumerate}
Causal interventions (Q1--Q4) are deliberately out of scope until a detector is validated.

\end{document}
